\section{Eastern Wisdom and Western Knowledge}

\textbf{by Ananda Coomaraswamy}

\emph{East and West}, \emph{The Crisis of the Modern World}, \emph{Introduction to the Study of} \emph{Hindu Doctrine}, and \emph{Man and His Becoming}, Luzac, London, (1941-1946) are the first in a series in which the majority of \textbf{Rene Guenon}'s works already published in French will appear in English. Another versions of \emph{Man and His Becoming} had appeared earlier. M. Rene Guenon is not an ``Orientalist'' but what the Hindus would call a ``master,'' formerly resident in Paris, and now for many years in Egypt, where his affiliations are Islamic. His \emph{Introduction general à l'étude des doctrines hindoues} appeared in 1921. As a preliminary to his further expositions of the traditional philosophy, sometimes called the \emph{Philosophia} \emph{Perennis (et Universalis)} must be understood, for this ``philosophy'' has been the common inheritance of all mankind without exception), Guenon cleared the ground of all possible misconceptions in two large and rather tedious, but by no means unnecessary, volumes, \emph{The Spiritist Fallacy} (a work for which \textbf{Bhagavad Gita}, XVII, 4, ``Men of darkness are they who make a cult of the departed and of spirits,'' might have served as a motto), and \emph{Theosophy, History of a Pseudo-Religion}.

These are followed by \emph{Man and his Becoming according to the Vedanta} and \emph{The Esoterism of Dante}, \emph{The King of the World}, \emph{St. Bernard}, \emph{East and West} and \emph{Spiritual Authority and Temporal Power}, \emph{The Symbolism of the Cross}, \emph{The Multiple States of Being}, and \emph{Oriental Metaphysics}.

More recently Guenon has published in mimeographed, and subsequently printed, editions \emph{The Reign of Quantity and the Sign of the Times}, and \emph{The Principles of Infinitesimal Calculus}.

In the meantime important articles from Guenon's pen appeared monthly in \emph{La Voile d'Isis}, later \emph{Études Traditionnelles}, a journal of which the appearance was interrupted by the war, but which has been continued as from September-October, 1945. \emph{Études Traditionnelles} is devoted to ``La Tradition Perpétuelle et Unanime, révélée tant par les dogmes et les rites des religions orthodoxies que par la langue universelle des symboles initiatiques''. The Perpetual and Unanimous Tradition, revealed as much by the dogmas and the rites of the orthodox religions as by the universal language of initiatory symbols. Of articles that have appeared elsewhere attention may be called to ``Islamic Esoterism'' in ``Cahiers du Sud''. Excerpts from Guenon's writings, with some comment, have appeared in `Triveni' (1935) and in the \emph{Visvabharati Quarterly} (1935,1938).

A work by L. De Gaigneron entitled \emph{Vers la connaissance interdite} is closely connected with Guenon's; it is presented in the form of a discussion in which the Atman (Spiritus), Mentalité (``Reason'', in the current, not the Platonic, sense), and a Roman abbé take part; the ``forbidden knowledge'' is that of the gnosis which the modern Church and the rationalist alike reject, though for different reasons --- the former because it cannot tolerate a point of view which considers Christianity only as one amongst other orthodox religions and the latter because, as a great Orientalist (Professor A. B. Keith) has remarked, ``such knowledge, as is not empirical, is meaningless to us and should not be described as knowledge'' --- an almost classical confession of the limitations of the ``scientific'' position. Guenon's French is at once precise and limpid, and inevitably loses in translation; his subject matter is of absorbing interest, at least to anyone who cares for what Plato calls the really serious things. Nevertheless it has often been found unpalatable; partly for reason already given, but also for reasons that have already been stated, paradoxically enough, by a reviewer of Blakney's Meister Eckhart in the `Harvard Divinity School Bulletin', who says that, ``To an age which believes in personality and personalism, the impersonality of mysticism is baffling; and to an age which is trying to quicken its insight into history the indifference of the mystics to events in time is disconcerting.'' As for history, Guenon's ``he who cannot escape from the standpoint of temporal succession so as to see all things in their simultaneity is incapable of the least conception of the metaphysical order'' adequately complements Jacob Boehme's designation of the ``history that was once brought to pass'' as ``merely the (outward) form of Christianity.'' For the Hindu, the events of the Rgveda are nowever and dateless, and the Krishna Lila ``not an historical event''; and the reliance of Christianity upon supposedly historical ``facts'' seems to be its greatest weakness.

The value of literary history for doxography is very little, and it is for this reason that so many orthodox Hindus have thought of Western scholarship as a ``crime'': their interest is not in ``what men have believed,'' but in the truth. A further difficulty is presented by Guenon's uncompromising language: ``Western civilization is an anomaly, not to say a monstrosity.'' Of this a reviewer has remarked that ``such sweeping remarks cannot be shared even by critics of Western achievements..''

I should have thought that now that its denouement is before our eyes, the truth of such a statement might have been recognized by every unprejudiced European; at any rate Sir George Birdwood in 1915 described modern Western civilization as ``Secular, joyless, inane, and self-destructive'' and Professor La Piana has said that ``What we call our civilization is but a murderous machine with no conscience and no ideals'' and might well have said suicidal as well as murderous. It would be very easy to cite innumerable criticisms of the same kind; Sir S. Radhakrishnan holds for example, that ``civilization is not worth saving if it continues on its present foundations,'' and this it would be hard to deny; Professor A. N. Whitehead has spoke quite as forcibly --- ``There remains the show of civilization, without any of it realities.''

In any case, if we are to read Guenon at all, we must have outgrown the temporally provincial view that has for so long and so complacently envisaged a continuous progress of humanity culminating in the twentieth century and be willing at least to ask ourselves whether there has not been rather a continued decline, ``from the stone until now,'' as one of the most learned men in the U.S.A. once put it to me. It is not by ``science'' that we can be saved:

\begin{quotex}
the possession of the sciences as a whole, if it does not include the best, will in some few cases aid but more often harm the owner.

We are obliged to admit that our European culture is a culture of the mind and senses only.

The prostitution of science may lead to world catastrophe.

Our dignity and our interests require that we shall be the directors and not the victims of technical and scientific advance.

Few will deny that the twentieth century thus far has brought us bitter disappointment.

We are now faced with the prospect of complete bankruptcy in every department of life.
\end{quotex}


Eric Gill speaks of the ``monstrous inhumanity'' of industrialism, and of the modern way of life, as ``neither human nor normal nor Christian... . It is our way of thinking that is odd and unnatural.'' This sense of frustration is perhaps the most encouraging sign of the times. We have laid stress on these things because it is only to those who feel this frustration, and not to those who still believe in progress, that Guenon addresses himself; to those who are complacent everything that he has to say will seem to be preposterous.

The reactions of Roman Catholics to Guenon are illuminating. One has pointed out that he is a ``serious metaphysician,'' i.e. one convinced of the truth he expounds and eager to show the unanimity of the Eastern and scholastic traditions, and observes that ``in such matters belief and understanding must go together.'' \emph{Crede ut intelligam} believe in order to understand is a piece of advice that modern scholars would, indeed, do well to consider; it is, perhaps, just because we have not believed that we have not yet understood the East. The same author writes of East and West,

\begin{quotex}
Rene Guenon is one of the few writers of our time whose work is really of importance ... he stands for the primacy of pure metaphysics over all other forms of knowledge, and presents himself as the exponent of a major tradition of thought, predominantly Eastern, but shared in the Middle Ages by the scholastics of the West ... clearly Guénon's position is not that of Christian orthodoxy, but many, perhaps most, of his theses are, in face, better in accord with authentic Thomist doctrine than are many opinions of devout but ill-instructed Christians.
\end{quotex}


We should do well to remember that even St. Tomas Aquinas did not disdain to make use of ``intrinsic and probable proofs'' derived form the ``pagan'' philosophers.

Gerald Vann, on the other hand, makes the mistake which the title of his review, ``Rene Guenon's Orientalism'' announces; for this is not another ``ism,'' nor a geographical antitheses but one of modern empiricism and traditional theory. Vann springs to the defense of the very Christianity in which Guenon himself sees almost the only possibility of salvation for the West; only possibility, not because there is no other body of truth, but because the mentality of the West is adapted to and needs a religion of just this sort. But if Christianity should fail, it is just because its intellectual aspects have been submerged, and it has become a code of ethics rather than a doctrine form which all other applications can and should be derived; hardly two consecutive sentences of some of Meister Eckhart's sermons would be intelligible to an average modern congregation, which does not expect doctrine, and only expects to be told how to behave.

If Guenon wants the West to turn to Eastern metaphysics, it is not because they are Eastern but because this is metaphysics. If ``Eastern'' metaphysics differed form a ``Western'' metaphysics --- as true philosophy differs from what is often so called in our modern universities --- one or the other would not be metaphysics. It is from metaphysics that the West has turned away in its desperate endeavor to live by bread alone, an endeavor of which the Dead Sea fruits are before our eyes. It is only because this metaphysics still survives as a living power in Eastern societies, in so far as they have not been corrupted by the withering touch of Western, or rather, modern civilization (for the contrast is not of East or West as such, but of ``those paths that the rest of mankind follows as a matter of course'' with those post-Renaissance paths that have brought us to our present impasse), and not to Orientalize the West, but to bring back the West to a consciousness of the roots of her own life and of values that have been transvalued in the most sinister sense, that Guenon asks us to turn to the East.

He does not mean, and makes it very clear that he does not mean, that Europeans ought to become Hindus or Buddhists, but much rather that they, who are getting nowhere by the study of ``the Bible as literature,'' or that of Dante ``as a poet,'' should rediscover Christianity, or what amounts to the same think, Plato (``that great priest,'' as Meister Eckhart calls him). I often marvel at men's immunity to the \emph{Apology}, \emph{Phaedo} or the seventh chapter of the \emph{Republic}; I suppose it is because they should not hear, ``thought one rose from the dead.''

The issue of ``East and West'' is not merely a theoretical (we must remind the modern reader that from the standpoint of the traditional philosophy, ``theoretical'' is anything but a term of disparagement) but also an urgent practical problem. Pearl Buck asks,

\begin{quotex}
Why should prejudices be so strong at this moment? The answer it seems to me is simple. Physical conveyance and other circumstances have forced parts of the world once remote from each other into actual intimacy for which peoples are not mentally or spiritually prepared. ... It is not necessary to believe that this initial stage must continue. If those prepared to act as interpreters will do their proper work, we may find that within another generation or two, or even sooner, dislike and prejudice may be gone. This is only possible if prompt and strong measures are taken by peoples to keep step mentally with the increasing closeness to which the war is compelling us.
\end{quotex}


But is this is to happen, the West will have to abandon what Guenon calls its ``proselytizing fury,'' an expression that must not be taken to refer only to the activities of Christian missionaries, regrettable as those often are, but of those of all the distributors of modern ``civilization'' and those of practically all those ``educators'' who feel that they have more to give than to learn form what are often called the ``backward'' or ``unprogressive'' peoples; to whom it does not occur that one may not wish or need to ``progress'' if one has reached a state of equilibrium that already provides for the realization of what one regards as the greatest purposed of life.

It is as an expression of good will and of the best intentions that this proselytizing fury takes on its most dangerous aspects. To many this ``fury'' can only suggest the fable of the fox that lost its tail, and persuaded the other foxes to cut off theirs. An industrialization of the East may be inevitable, but do not let us call it a blessing that a folk should be reduced to the level of a proletariat, or assume that materially higher standards of living necessarily make for greater happiness. The West is only just discovering, to its great astonishment, that ``material inducements, that is, money or the things that money can buy'' are by no means so cogent a force as has been supposed; ``Beyond the subsistence level, the theory that this incentive is decisive is largely an illusion.''

As for the East, as Guenon says, ``The only impression that, for example, mechanical invention make on most Orientals is one of deep repulsion; certainly it all seems to them far more harmful than beneficial, and if they find themselves obliged to accept certain things which the present epoch has made necessary, they do so in the hope of future riddance ... what the people of the West call `rising' would be called some `sinking'; that is what all true Orientals think.'' It must not be supposed that because so many Eastern peoples have imitated us in self-defense that they have therefore accepted our values; on the contrary, it is just because the conservative East still challenges all the presuppositions on which our illusion of progress rests, that is deserves our most serious consideration.

There is nothing in economic intimacies that is likely to reduce prejudice or promote mutual understanding automatically. Even when Europeans live amongst Orientals, ``economic contact between the Eastern and Western groups is practically the only contact there is. There is very little social or religious give and take between the two. Each lives in a world almost entirely closed to the other --- and by `closed' we man not only `unknown' but more: incomprehensible and unattainable.'' That is an inhuman relationship, by which both parties are degraded.

Neither must it be assumed that the Orient thinks it important that the masses should learn to read and write. Literacy is a practical necessity in an industrial society, where the keeping of account is all important. But in India, in so far as Western methods of education have not been imposed from without, all higher education is imparted orally, and to have heard is for more important than to have read. At the same time the peasant, prevented by his illiteracy and poverty from devouring the newspapers and magazines that form the daily and almost the only reading of the vast majority of Western ``literates,'' is, like Hesiod's Boeotian farmers, and still more like the Gaelic-speaking Highlanders before the era of the board schools, thoroughly familiar with an epic literature of profound spiritual significance and a body of poetry and music of incalculable value; and one can only regret the spread of an ``education'' that involves the destruction of all these things, or only preserves them as curiosities within the covers of books.

For cultural purposed it is not important that the masses should be literate; it is not necessary that anyone should be literate; it is only necessary that there should be amongst the people philosophers (in the traditional, not the modern sense of the word), and that there should be preserved deep respect on the part of laymen for true learning that is the antithesis of the American attitude to a ``professor''. In these respects the whole East is still far in advance of the West, and hence the learning of the elite exerts a far profounder influence upon society as a whole than the Western specialist ``thinker'' can ever hope to wield.

It is not, however, primarily with a protection of the East against the subversive inroads of Western ``culture'' that Guenon is concerned, but rather with the question, What possibility of regeneration, if any, can be envisaged for the West? The possibility exists only in the event of a return to first principles and to the normal ways of living that proceed from the application of first principles to contingent circumstances; and as it is only in the East that these things are still alive, it is to the East that the West must turn. ``It is the West that must take the initiative, but she must be prepared really to go towards the East, not merely seeking to draw the East towards herself, as she has tried to do so far. There is no reason why the East should take this initiative, and there would still be none, even if the Western world were not in such a state as to make any effort in this direction useless. ... It now remains for us to show how the West might attempt to approach the East.''

He proceeds to show that the work is to be done in the two fields of metaphysics and religion, and that it can only be carried out on the highest intellectual levels, where agreement on first principles can be reached apart from any propaganda on behalf of or even apology for ``Western civilization.''

The work must be undertaken, therefore, by an ``elite.'' And as it is here more than anywhere that Guenon's meaning is likely to be willfully misinterpreted, we must understand clearly what he manes by such an elite. The divergence of the West and East being only ``accidental,'' ``the bringing of these two portions of mankind together and the return of the West to a normal civilization are really just one and the same thing.'' An elite will necessarily work in the first place ``for itself, since its members will naturally reap from their own development an immediate and altogether unfailing benefit.'' An indirect result --- ``indirect,'' because on this intellectual level one does not think of ``doing good'' to others, or in terms of ``service,'' but seeks truth because one needs it oneself --- would, or might under favorable conditions, bring about ``a return of the West to a traditional civilization,'' i.e. one in which ``everything is seen as the application and extension of a doctrine whose essence is purely intellectual and metaphysical.''

It is emphasized again and again that such an elite does not mean a body of specialists or scholars who would absorb and put over on the West the forms of an alien culture, nor even persuade the West to return to such a traditional civilization as existed in the Middle Ages. Traditional cultures develop by the application of principles to conditions; the principles, indeed, are unchangeable and universal, but just as \textbf{nothing can be known except in the mode of the knower}, so nothing valid can be accomplished socially without taking into account the character of those concerned and the particular circumstances of the period in which they live.

There is no ``fusion'' of cultures to be hoped for; it would be noting like an ``eclecticism'' or ``syncretism'' that an elite would have in view. Neither would such an elite be organized in any way so as to exercise such a direct influence as that which, for example, the Technocrats would like to exercise for the good of mankind. If such an elite ever came into being, the vast majority of Western men would never know of it; it would operate only as a sort of leaven, and certainly on behalf of rather than against whatever survives of traditional essence in, for example the Greek Orthodox and Roman Catholic domains. It is, indeed, a curious fact that some of the most powerful defenders of Christian dogma are to be found amongst Orientals who are not themselves Christians, or ever likely to become Christians, but recognize in the Christian tradition an embodiment of the universal truth to which God has never nor anywhere left himself without a witness.

In the meantime, Guenon asks, ``Is this really `the beginning of an end' for the modern civilization? ... At least there are many signs which should give food for reflection to those who are still capable of it; will the West be able to regain control of herself in time?'' Few would deny that we are faced with the possibility of a total disintegration of culture. We are at war with ourselves, and therefore at war with one another. Western man is unbalanced, and the question, Can he recover himself? is a very real one. No one to whom the question presents itself can afford to ignore the writing of the leading living exponent of a traditional wisdom that is no more essentially Oriental than it is Occidental, though it may be only in the uttermost parts of the earth that it is still remembered and must be sought.


\flrightit{Posted on 2012-10-27 by Cologero}

